%%%%%%%%%%%%%%%%%%%%%part6.tex%%%%%%%%%%%%%%%%%%%%%%%%%%%%%%%%%%
% 
% Part VI: Verifiable Architecture & Autonomous Agents
%
%%%%%%%%%%%%%%%%%%%%%%%% Springer Nature %%%%%%%%%%%%%%%%%%%%%%

\begin{partbacktext}
\part{Verifiable Architecture \& Autonomous Agents}
\noindent In this sixth part of the book, we move beyond the model to the broader system architecture. We explore how to build AI systems that are not just trustworthy by design, but \textit{verifiably} trustworthy through cryptography and formal logic.

We start in \textbf{Chapter 21} with \textbf{Zero-Knowledge Machine Learning (ZK-ML)}, providing a mechanism for verifying model execution without revealing private data or weights. \textbf{Chapter 22} introduces \textbf{Formal Verification}, applying rigorous mathematical proof to ensure model safety. Finally, in \textbf{Chapter 23}, we address the emerging challenge of \textbf{Agentic AI Governance}, ensuring that autonomous agents respect privacy boundaries and act as intended.

\begin{casestudy} 
The 2025 collapse of VeriBank"s external auditing process occurred when it was discovered that human auditors had missed over 40,000 fraudulent transactions hidden by a sophisticated adversarial ML model. The realization that human oversight alone could not keep pace with machine-speed fraud led to the "Verification Turn" of 2026, making the Zero-Knowledge and Formal Verification methods in Part VI an architectural requirement. 
\end{casestudy} 
\end{partbacktext}
